\chapter{Matching versus Parsing}
\label{sec:matching-vs-parsing}

Chapter~\ref{sec:derivatives} established matching with a Boolean
answer: whether $w \in L(r)$, the word is either accepted or rejected.
Parsing refines this question by asking for the derivation that
witnesses the match. Parse trees provide a representation of such
derivations, recording the choices made by the structure of $r$ while
consuming $w$.

This refinement introduces ambiguity: a fixed $r$ and $w$ may admit
multiple parse trees. A parser must therefore combine the semantics of
matching with a \emph{disambiguation policy} that selects one of them.
This chapter introduces parse trees, examines Greedy and POSIX
disambiguation, and establishes the ordering that the parsers developed
later must compute.

%

\section{The Matching Problem}
\label{sec:matching-problem}

The three matchers of Chapter~\ref{sec:derivatives} all answer the same
Boolean question: given $r$ and $w$, is $w \in L(r)$? Several
applications need more than that answer. For example
knowing \emph{which alternative} of $r$ matched, and
knowing \emph{which part} of $w$ a given subexpression matched. Boolean
recognition discards both kinds of information: it confirms that some
derivation of the match exists, but nothing about what that derivation
actually was \citep{frischcardelli2004}.

\paragraph{Which alternative: lexing and schema validation.}
Consider a lexer for a programming language, using the pattern
\texttt{if|[a-z]+} to recognize a single token that could be the
keyword \texttt{if}, marking the start of a conditional statement, or
an ordinary lowercase identifier, naming a variable or function.
Against the input \texttt{"if"}, both readings apply: the string is the
literal keyword, but it is also a valid identifier, since it consists
only of lowercase letters. A Boolean matcher confirms only that the
pattern matches \texttt{"if"}; it does not say which of the two
readings was taken, and the two readings are not interchangeable
downstream. In a statement such as \texttt{if (x > 0) foo();}, the
parser needs an \texttt{IF} token to recognize a conditional. If the
lexer instead reported an identifier, \texttt{IDENT("if")}, the parser
would try to read \texttt{if} as an expression, and the conditional
would never be recognized.

In practice, this particular ambiguity is resolved before parsing ever
begins, not during it. Most languages forbid identifiers from being
spelled the same as a keyword, and lexer generators enforce this
directly: flex always matches the longest possible token, and when two
rules match input of the same length, the rule listed earlier in the
specification wins \citep{flex-manual-patterns}. Placing the keyword
pattern for \texttt{if} before the general identifier pattern is
therefore enough to guarantee the keyword reading whenever both would
otherwise apply, and \texttt{"if"} can never be tokenized as an
identifier at all. The underlying problem, that a Boolean matcher
cannot report which alternative explains a match, has not gone away; it
has been settled in advance by a rule the tool applies uniformly,
rather than recovered from any particular match. That rule is already a
disambiguation policy in miniature, of exactly the kind
Section~\ref{sec:greedy-leftmost} and
Section~\ref{sec:posix-longest-leftmost} formalize later in this
chapter.

Schema validation avoids the same problem whenever a schema declares a
field's type up front: checking a value declared as an integer against
the pattern for integers never raises the question of whether it might
also be a Boolean, since no other pattern is ever consulted. The
problem reappears in formats that infer a value's type from its text
instead of from a declared schema. YAML is the clearest documented case.
A plain, unquoted scalar is resolved to a type by checking it against an
ordered set of patterns, one for null, one for booleans, one for
integers, one for floats, and so on, with the type belonging to the
first pattern that matches \citep{yaml-1-1-spec}. In YAML 1.1, the
boolean pattern accepts not only \texttt{true} and \texttt{false} but
also \texttt{yes}, \texttt{no}, \texttt{on}, \texttt{off}, and several
capitalizations of each, so the two-letter ISO country code for Norway,
\texttt{NO}, is resolved to the Boolean value \texttt{false} rather than
the string \texttt{"NO"} in any YAML document that does not quote it
explicitly. This became widely known as the "Norway problem," and it is
exactly the failure this section describes: the text \texttt{NO}
matches both the boolean pattern and, trivially, the fallback string
pattern, and only the fixed order in which YAML checks its candidate
types, not anything about the text \texttt{NO} itself, decides which one
wins \citep{strictyaml-implicit-typing}. YAML 1.2 removed the
\texttt{yes}/\texttt{no}/\texttt{on}/\texttt{off} boolean forms
specifically to shrink this overlap, though widely used implementations
such as PyYAML still default to the 1.1 rules.

%

\paragraph{A pattern where the two readings genuinely compete.}
The \texttt{"if"} example is resolved by a tie-breaking convention. Both
alternatives of \texttt{if|[a-z]+} match the same two-character input, so
a longest-match rule cannot distinguish them. The lexical specification
therefore gives the keyword alternative priority. What appears to be an
obvious choice is thus the result of a specific disambiguation rule: the
two alternatives consume equally long substrings, and the earlier
alternative wins the tie.

The pattern $(a+ab)(b+\varepsilon)$ against \texttt{"ab"} presents a
different situation. The two alternatives of $a+ab$ consume different
amounts of input. The \texttt{ab} alternative can consume both characters,
after which $(b+\varepsilon)$ matches $\varepsilon$. Alternatively, the
\texttt{a} alternative consumes only the first character, leaving
\texttt{b} to be matched by the \texttt{b} alternative of
$(b+\varepsilon)$. Both choices produce a complete and valid match of the
same input.

Here, two natural disambiguation policies give different results. A
policy that prefers the longer match selects the \texttt{ab} alternative,
whereas a policy that gives priority to the earlier alternative selects
\texttt{a}. Unlike the \texttt{"if"} example, neither choice can be
dismissed as merely resolving a tie. The expression admits both parses,
and selecting between them requires an explicit policy.

Before such choices can be compared, however, the possible parses must be
made explicit. A Boolean result records only that a match exists; it does
not record the structural choices through which the match was obtained.
Section~\ref{sec:parse-trees-as-evidence} represents these choices as
\emph{parse trees} for a structured \emph{witness} to a
match: it records how the structure of $r$ was used to consume $w$. When
multiple parse trees witness the same match, the parser must determine
which one to return.

%

\section{Parse Trees as Matching Evidence}
\label{sec:parse-trees-as-evidence}

Sulzmann and Lu formalize such a witness by viewing regular expressions as
types and parse trees as the values inhabiting them
\citep{sulzmannlu2014}. A parse tree $v$ is built from the same
constructors used informally in Section~\ref{sec:regex-semantics}'s
definition of $L(r)$: a pair for concatenation, a tagged left or right
injection for alternation, and a list for repetition. A judgment
$\vdash v : r$, defined by structural recursion on both $v$ and $r$,
states that $v$ is a valid parse tree for $r$, by the rules of
Equation~\eqref{eq:typing-rules}:

\begin{equation}\label{eq:typing-rules}
\begin{array}{ll}
\vdash () : \varepsilon \quad (\mathsf{T}_{\varepsilon}) &
\dfrac{l \in \Sigma}{\vdash l : l} \quad (\mathsf{T}_{\mathrm{lit}}) \\[1.2em]
\dfrac{\vdash v_1 : r_1 \quad \vdash v_2 : r_2}
      {\vdash (v_1,v_2) : r_1 r_2} \quad (\mathsf{T}_{\mathrm{seq}})
&
\dfrac{\vdash v_1 : r_1}
      {\vdash \mathit{Left}\ v_1 : r_1 + r_2} \quad (\mathsf{T}_{+}^{L})
\quad
\dfrac{\vdash v_2 : r_2}
      {\vdash \mathit{Right}\ v_2 : r_1 + r_2} \quad (\mathsf{T}_{+}^{R})
\\[1.2em]
\dfrac{}{\vdash [\,] : r^{*}} \quad (\mathsf{T}_{*}^{0}) &
\dfrac{\vdash v : r \quad \vdash vs : r^{*}}
      {\vdash (v:vs) : r^{*}} \quad (\mathsf{T}_{*})
\end{array}
\end{equation}

As with the derivative rules of
Section~\ref{sec:brzozowski-derivatives}, each rule carries a name, and
the derivations in this chapter cite those names together with the
binding of each metavariable, so that it is always visible which
subexpression of $r$ a given $r_1$ or $r_2$ refers to.

A \emph{flattening} function $\lvert v\rvert$ reads the word represented
by a parse tree by structural recursion:

\begin{equation}\label{eq:flatten}
\begin{aligned}
\lvert ()\rvert &= \varepsilon, &
\lvert l\rvert &= l, &
\lvert [\,]\rvert &= \varepsilon, \\
\lvert (v_1,v_2)\rvert &= \lvert v_1\rvert\lvert v_2\rvert, &
\lvert \mathit{Left}\ v\rvert &= \lvert v\rvert, &
\lvert \mathit{Right}\ v\rvert &= \lvert v\rvert, \\
\lvert v:vs\rvert &= \lvert v\rvert\lvert vs\rvert. &&
\end{aligned}
\end{equation}

The key property is that the parse trees characterize the language:

\begin{equation}\label{eq:mvp-5}
L(r) = \{\,\lvert v\rvert \mid \vdash v : r\,\}.
\end{equation}

Thus, a parse tree together with its derivation provides a structured
\emph{witness} that $w \in L(r)$: it records not only that the match
exists, but the choices made by the regular expression while producing
$w$.

%


\paragraph{Unambiguous example.}
For an unambiguous match, this witness is unique. For
$r = a \cdot b$ and $w = \mathtt{ab}$, only one derivation is possible.
Matching $r$ against $\mathsf{T}_{\mathrm{seq}}$'s conclusion binds
$r_1 = a$ and $r_2 = b$, so the rule requires $\vdash v_1 : a$ and
$\vdash v_2 : b$. Each of those is a literal, and
$\mathsf{T}_{\mathrm{lit}}$ is the only rule whose conclusion matches a
literal, so $v_1 = a$ and $v_2 = b$ with no choice anywhere:

\begin{equation}\label{eq:mvp-6}
\dfrac{\dfrac{a \in \Sigma}{\vdash a : a}\ (\mathsf{T}_{\mathrm{lit}})
 \quad
 \dfrac{b \in \Sigma}{\vdash b : b}\ (\mathsf{T}_{\mathrm{lit}})}
      {\vdash (a,b) : ab}\ (\mathsf{T}_{\mathrm{seq}})
\end{equation}

This is the only possible derivation, yielding $v = (a,b)$, and
flattening it recovers the input:

\begin{equation}\label{eq:mvp-1}
\begin{aligned}
\lvert (a,b)\rvert
  &= \lvert a\rvert\,\lvert b\rvert
  &&\text{flattening of a pair,}\ v_1 = a,\ v_2 = b \\
  &= \mathtt{ab}
  &&\text{flattening of a literal, twice}
\end{aligned}
\end{equation}

When the structure of the expression permits different choices, the
witness itself can be ambiguous. The first two examples below are
adapted from the running examples of \citet{sulzmannlu2014} and are
reproduced in this project's test suite as \texttt{paper\_r1} and
\texttt{paper\_r2}. They are included here because the same cases are
used later to illustrate and evaluate the parser implementations
developed in Chapters~\ref{sec:derivative-parser} and
\ref{sec:partial-derivative-parser}. The derivations given here provide
a reference for that later discussion. Two further examples, drawn from
this project's own test suite, follow.

%

\paragraph{Ambiguous example 1: concatenation.}
Consider $r_1 = (a + ab)(b + \varepsilon)$ against $w = \mathtt{ab}$.
One valid derivation chooses the second alternative of $a+ab$, which
consumes both characters, and the empty alternative of
$b+\varepsilon$:

\begin{equation}\label{eq:mvp-7}
\dfrac{
  \dfrac{\vdash a : a \quad \vdash b : b}
        {\vdash (a,b) : ab}
}{
  \vdash \mathit{Right}(a,b) : a + ab
}
\qquad
\dfrac{
  \vdash () : \varepsilon
}{
  \vdash \mathit{Right}() : b + \varepsilon
}
\end{equation}

and therefore

\begin{equation}\label{eq:mvp-8}
\vdash
\bigl(\mathit{Right}(a,b),\mathit{Right}()\bigr)
:
(a+ab)(b+\varepsilon).
\end{equation}

The corresponding parse tree is

\begin{equation}\label{eq:mvp-9}
v
=
\bigl(\mathit{Right}(a,b),\mathit{Right}()\bigr),
\end{equation}

with

\begin{equation}\label{eq:mvp-10}
\lvert v\rvert
=
\lvert (a,b)\rvert\,\lvert ()\rvert
=
\mathtt{ab}\,\varepsilon
=
\mathtt{ab}.
\end{equation}

This is not the only witness. The first alternative of $a+ab$ can
instead consume $\mathtt{a}$, leaving $\mathtt{b}$ for the first
alternative of $b+\varepsilon$:

\begin{equation}\label{eq:mvp-11}
\dfrac{
  \vdash a : a
}{
  \vdash \mathit{Left}\ a : a + ab
}
\qquad
\dfrac{
  \vdash b : b
}{
  \vdash \mathit{Left}\ b : b + \varepsilon
}
\end{equation}

and therefore

\begin{equation}\label{eq:mvp-12}
\vdash
\bigl(\mathit{Left}\ a,\mathit{Left}\ b\bigr)
:
(a+ab)(b+\varepsilon).
\end{equation}

The corresponding parse tree is

\begin{equation}\label{eq:mvp-13}
v'
=
\bigl(\mathit{Left}\ a,\mathit{Left}\ b\bigr),
\end{equation}

with

\begin{equation}\label{eq:mvp-14}
\lvert v'\rvert
=
\lvert a\rvert\,\lvert b\rvert
=
\mathtt{ab}.
\end{equation}

Thus the same expression $r_1$ and word $w$ have two distinct parse
trees, $v$ and $v'$, satisfying

\begin{equation}\label{eq:mvp-15}
\lvert v\rvert = \lvert v'\rvert = \mathtt{ab}.
\end{equation}

%

\paragraph{Ambiguous example 2: repetition.}
The same phenomenon can arise from repetition rather than
concatenation. Consider $r_2 = (a+b+ab)^{*}$ against $w = \mathtt{ab}$,
nesting the alternation as $a+(b+ab)$
(\texttt{paper\_r2}\footnote{The associativity of $+$ is immaterial to
$L(r)$ and to this example.} in the test suite). One derivation uses a
single iteration, taking the \texttt{ab} alternative directly:

\begin{equation}\label{eq:mvp-16}
\dfrac{
  \dfrac{
    \dfrac{\vdash a : a \quad \vdash b : b}
          {\vdash (a,b) : ab}
  }{
    \vdash \mathit{Right}(a,b) : b + ab
  }
}{
  \vdash \mathit{Right}\ \mathit{Right}(a,b) : a + (b+ab)
}
\end{equation}

which becomes a one-element list via the star's constructor rule:

\begin{equation}\label{eq:mvp-17}
\vdash [\mathit{Right}\ \mathit{Right}(a,b)] : (a+(b+ab))^{*}.
\end{equation}

The corresponding parse tree is

\begin{equation}\label{eq:mvp-18}
v_2 = [\mathit{Right}\ \mathit{Right}(a,b)],
\qquad
\lvert v_2\rvert = \lvert (a,b)\rvert = \mathtt{ab}.
\end{equation}

A second derivation instead uses two iterations, matching $\mathtt{a}$
and $\mathtt{b}$ through separate alternatives of the star's body:

\begin{equation}\label{eq:mvp-19}
\dfrac{\vdash a : a}
      {\vdash \mathit{Left}\ a : a+(b+ab)}
\qquad
\dfrac{
  \dfrac{\vdash b : b}
        {\vdash \mathit{Left}\ b : b+ab}
}{
  \vdash \mathit{Right}\ \mathit{Left}\ b : a+(b+ab)
}
\end{equation}

which combine, one iteration at a time, into

\begin{equation}\label{eq:mvp-20}
\vdash [\mathit{Left}\ a,\ \mathit{Right}\ \mathit{Left}\ b] :
(a+(b+ab))^{*}.
\end{equation}

The corresponding parse tree, with two iterations of the star, is

\begin{equation}\label{eq:mvp-21}
v_2' = [\mathit{Left}\ a,\ \mathit{Right}\ \mathit{Left}\ b],
\qquad
\lvert v_2'\rvert = \lvert a\rvert\,\lvert b\rvert = \mathtt{ab}.
\end{equation}


As with $r_1$, both $v_2$ and $v_2'$ witness the same match: ambiguity
here comes from how many iterations the star takes, not from which
top-level alternative fires.

%

\paragraph{Ambiguous example 3: duplicate alternatives.}
Concatenation and repetition are not the only sources of ambiguity.
Consider $r_3 = (a+a)^{*}$ against $w = \mathtt{aa}$. Here the star's
body contains two syntactically distinct but semantically identical
alternatives for the same literal. This is an instance of
\texttt{ambiguous\_star\_dedup\_matches\_bitcoded} in
\texttt{src/parsers/pderiv\_std/parse.rs}
\footnote{The test uses a four-character input; the two-character case
shown here exhibits the same phenomenon.}.

Each of the two positions in $w$ can
independently be matched by either the left or the right copy of
\texttt{a}, giving four distinct parse trees for this two-character
input alone:

\begin{equation}\label{eq:mvp-22}
[\mathit{Left}\ a,\mathit{Left}\ a],\quad
[\mathit{Left}\ a,\mathit{Right}\ a],\quad
[\mathit{Right}\ a,\mathit{Left}\ a],\quad
[\mathit{Right}\ a,\mathit{Right}\ a],
\end{equation}

all flattening to $\mathtt{aa}$. Nothing about $\mathit{Left}\ a$ and
$\mathit{Right}\ a$ distinguishes them semantically; the ambiguity is
entirely an artifact of writing the same alternative twice. This kind
of duplication is particularly relevant to the partial-derivative
parser, where distinct syntactic paths can produce equivalent residual
matches. The deduplication mechanisms discussed in
Chapter~\ref{sec:partial-derivative-parser} prevent such redundant
paths from accumulating in the parse frontier.


%

\paragraph{Ambiguous example 4: an IPv4 pattern.}
The previous examples isolate particular sources of ambiguity. The same
phenomenon also occurs in a pattern taken from practice: a regular
expression for dotted-quad IPv4 addresses from the Snort
intrusion-detection dataset studied by \citet{mamouras2024}. The pattern
is included in this project's test suite as
\texttt{snort\_ip\_regex}\footnote{\texttt{tests/test\_frontend.rs}.}.

\begin{center}
\texttt{((\textbackslash d[1-9]\textbackslash d|1\textbackslash d\textbackslash d|2[0-4]\textbackslash d|25[0-5])\textbackslash .)\{3\}}\\
\texttt{(\textbackslash d[1-9]\textbackslash d|1\textbackslash d\textbackslash d|2[0-4]\textbackslash d|25[0-5])}
\end{center}

using the class and repetition shorthands of
Section~\ref{sec:frontend-translate}, under which \texttt{\textbackslash
d} desugars to a ten-way alternation of the literals \texttt{0}
through \texttt{9} and \texttt{\{3\}} desugars to three concatenated
copies of its operand. The pattern concatenates four copies of a
shared three-character octet pattern, separated by literal dots. Write
$O = A+B+C+D$ for this shared four-way alternation, with

\begin{center}
$A = $ \texttt{\textbackslash d[1-9]\textbackslash d} \quad
$B = $ \texttt{1\textbackslash d\textbackslash d} \quad
$C = $ \texttt{2[0-4]\textbackslash d} \quad
$D = $ \texttt{25[0-5]},
\end{center}

left-associated as $((A+B)+C)+D$, so that the full pattern is
$(O\,\mathtt{.})^{3}O$ with the final dot omitted. Already at the level
of a single octet, $O$ is ambiguous on some inputs in exactly the sense
of $r_1$: against $w=\mathtt{239}$, alternative $A$ matches, reading
the digits as \emph{any digit}, \emph{a digit from \texttt{1} to
\texttt{9}}, \emph{any digit}; but so does alternative $C$, reading
\texttt{2} as the literal \texttt{2}, \texttt{3} as a digit from
\texttt{0} to \texttt{4}, and \texttt{9} as any digit ($B$ and $D$ both
fail, since $w$ does not begin with \texttt{1} or \texttt{25}). 
Both $A$ and $C$ consume all three characters of $w$. Thus, unlike
$r_3$, where the alternatives are semantically identical, the ambiguity
here is between two genuinely different alternatives that happen to
match the same substring. \citet{mamouras2024} use this same pattern to
illustrate a different symptom of the same underlying policy
difference: under unanchored substring search, Greedy's leftmost-first
matching can return a shorter, truncated address where POSIX returns
the complete one, a distinction in match \emph{span} rather than, as
here, in which of two total-match parse trees is selected.


%


\section{Disambiguation Policies}
\label{sec:disambiguation-problem}

The parse-tree judgment identifies every valid witness, but it does not
say which witness a parser should return. A \emph{disambiguation policy}
adds this missing choice. Given a regular expression $r$ and word $w$, the
policy selects one parse tree $v$ satisfying

\begin{equation}\label{eq:mvp-23}
\vdash v : r
\qquad\text{and}\qquad
\lvert v\rvert = w.
\end{equation}

The four ambiguous examples above illustrate why the policy matters, and
in more than one way. For $r_1$ and $\mathtt{ab}$, one witness chooses
$\texttt{a}$ before $\texttt{b}$, while the other assigns the whole input
to $\texttt{ab}$. For $r_2$ and $\mathtt{ab}$, one witness uses two
iterations and the other uses one. For $r_3$ and $\mathtt{aa}$, four
witnesses differ only in which of two identical alternatives each
position is attributed to. For $O$ and $\mathtt{239}$, two alternatives
with different content, not merely different labels, both happen to
consume the whole octet. Different matching conventions select
different witnesses, and, as $r_3$ and $O$ show, a policy also has to
say something about witnesses that agree on every character consumed
and differ only in which alternative gets the credit.

The two policies considered here resolve these choices differently.
Greedy matching follows the structural priority of the expression,
whereas POSIX matching gives priority to longer leftmost submatches.
Their difference becomes clear when they are applied to the same
ambiguous expression.


%


\subsection{Greedy Leftmost Matching}
\label{sec:greedy-leftmost}

The policy associated with conventional backtracking matchers is
\emph{Greedy}. Alternatives are considered from left to right, and the
first alternative that leads to an overall match is selected. 
Within a Kleene star, another iteration is preferred to stopping whenever
continuing the repetition can still lead to an overall match
\citep{frischcardelli2004}.

For alternatives, the defining preference is structural: the left
alternative is preferred regardless of how much input the two
alternatives consume. In an ordering notation, the essential rule is

\begin{equation}\label{eq:mvp-24}
\mathit{Left}\ v
\succ_{r_1+r_2}
\mathit{Right}\ v'.
\end{equation}

Every other comparison rule: (C1) and (C2) for concatenation, (A3)
and (A4) for continuing a comparison once both alternatives already
agree on a branch, and (K1)--(K4) for repetition is shared,
unchanged, between Greedy and POSIX; the two orderings differ only in
how they compare two alternatives that choose \emph{different}
branches, where POSIX replaces the rule above with the length-comparing
pair (A1) and (A2) of Section~\ref{sec:posix-longest-leftmost}.
Sulzmann and Lu construct their
POSIX ordering explicitly as this one substitution into Frisch and
Cardelli's Greedy ordering \citep{frischcardelli2004,sulzmannlu2014},
leaving every other rule untouched.

%

\paragraph{Applied to $r_1$.}
Greedy first tries the \texttt{a} alternative. It consumes the leading
\texttt{a}, and the remaining \texttt{b} is successfully matched by the
first alternative of $b+\varepsilon$. Since the complete match succeeds,
there is no need to try the \texttt{ab} alternative. The Greedy result is
therefore

\begin{equation}\label{eq:mvp-25}
v'
=
\bigl(\mathit{Left}\ a,\mathit{Left}\ b\bigr).
\end{equation}

\paragraph{Applied to $r_2$.}
Greedy similarly begins with the first alternative, matching $\mathtt{a}$.
The star then attempts another iteration, in which the first alternative
fails on $\mathtt{b}$ and the second alternative succeeds. The resulting
parse uses two iterations:

\begin{equation}\label{eq:mvp-26}
v_2'
=
[\mathit{Left}\ a,\mathit{Right}\ \mathit{Left}\ b].
\end{equation}

The Greedy result therefore depends directly on the order in which the
regular expression presents its alternatives.

\paragraph{Applied to $r_3$.}
At each of the two positions, $\mathit{Left}\ v \succ_{a+a} \mathit{Right}\
v'$ applies directly, regardless of length, since both alternatives
consume the same single character. Greedy therefore selects the left
alternative every time:

\begin{equation}\label{eq:mvp-27}
[\mathit{Left}\ a,\mathit{Left}\ a].
\end{equation}

Unlike $r_1$ and $r_2$, this result is not decided by which alternative
matches more of the input; every alternative here matches exactly one
character. It is decided purely by position in the pattern, the same
kind of tie-breaking that Section~\ref{sec:matching-problem} showed
resolving \texttt{"if"} against \texttt{if|[a-z]+}.

\paragraph{Applied to $O$.}
Against $w=\mathtt{239}$, Greedy tries $O$'s alternatives left to
right: $A$, then $B$, then $C$, then $D$. Alternative $A$ matches the
leading digit \texttt{2} against \texttt{\textbackslash d}, the digit
\texttt{3} against \texttt{[1-9]}, and the digit \texttt{9} against
\texttt{\textbackslash d}, consuming all three characters. Since this
already produces a complete match, Greedy never tries $B$, $C$, or
$D$, even though $C$ would also have matched. The Greedy result
selects $A$, exactly as it selected the left alternative of $r_1$ and
of $r_3$: once an alternative leads to a full match, later
alternatives are never even attempted.

%

\paragraph{Procedural versus declarative view.}
The description above gives a procedural intuition for Greedy matching:
alternatives are considered from left to right, and a choice is
retained when it can be extended to a complete match. The relation
$\succ_{r_1+r_2}$ used throughout this subsection is a different kind
of thing: a declarative ordering on the parse trees admitted by $r$
and $w$, generalized to arbitrary $r$ as $\succ_r$ in
Section~\ref{sec:posix-longest-leftmost}. \citet{frischcardelli2004},
whose ordering this one adapts, are explicit that these two views are
not the same thing computationally: a backtracking procedure computing
their ordering's maximum is possible, but they describe it as ``quite
inefficient because of backtracking,'' and spend the rest of their
paper building a two-pass, backtracking-free algorithm instead, with a
proof that it computes the same value. Neither \texttt{pderiv\_std} nor
\texttt{pderiv\_bc} (Chapter~\ref{sec:partial-derivative-parser})
backtracks either, for the same reason: both compute the Greedy answer
through a forward partial-derivative frontier with
first-occurrence-wins deduplication. The distinction matters
concretely in Section~\ref{sec:infinite-empty-chain}, where the
ordering can have no maximal element at all, even though a concrete,
non-backtracking parser can still compute the intended Greedy result.

%

\subsection{The Infinite Empty-Match Chain Problem}
\label{sec:infinite-empty-chain}

Greedy's preference for continuing a repetition creates a problem when
the repeated expression can match the empty word. Consider
$\varepsilon^{*}$ against the empty input. The following parse trees are
all valid:

\begin{equation}\label{eq:mvp-28}
v_0 = [\,],
\qquad
v_1 = [()],
\qquad
v_2 = [(),()],
\qquad \ldots
\end{equation}

Each additional iteration still produces the same flattened word
$\varepsilon$. Under the Greedy ordering, every additional iteration is
preferred:

\begin{equation}\label{eq:mvp-29}
v_0
\prec_{\varepsilon^{*}}
v_1
\prec_{\varepsilon^{*}}
v_2
\prec_{\varepsilon^{*}}
\cdots.
\end{equation}

There is therefore no maximal parse tree. The problem is not that the
language is ill-defined; $\varepsilon$ is clearly a valid match. The
problem is that the Greedy preference does not identify a largest witness
when empty iterations can be added indefinitely.

A slightly less degenerate example is
$(\varepsilon+a)^{*}$ against $\mathtt{a}$. The parse trees

\begin{equation}\label{eq:mvp-30}
\begin{aligned}
v_0 &= [\mathit{Right}\ a],\\
v_1 &= [\mathit{Left}\ (),\mathit{Right}\ a],\\
v_2 &= [\mathit{Left}\ (),\mathit{Left}\ (),\mathit{Right}\ a],\\
&\quad\vdots
\end{aligned}
\end{equation}

all flatten to $\mathtt{a}$. Greedy prefers each additional empty
iteration, producing another infinite ascending chain.

The same pathology reappears whenever a star's body is itself nullable,
not just when the body is $\varepsilon$ outright. Take $(a^{*})^{*}$
against $\mathtt{aaa}$: since the inner $a^{*}$ can always match
$\varepsilon$, an outer iteration contributing nothing is always
available, and Greedy always prefers taking it:

\begin{equation}\label{eq:mvp-31}
[[a,a,a]]
\;\prec_{(a^{*})^{*}}\;
[[a,a,a],[\,]]
\;\prec_{(a^{*})^{*}}\;
[[a,a,a],[\,],[\,]]
\;\prec_{(a^{*})^{*}}\;
\cdots,
\end{equation}

another infinite ascending chain, this time built on top of an
otherwise unambiguous inner match rather than on $\varepsilon$ itself.
Nested, mutually nullable stars of this shape are also the classic
trigger for catastrophic backtracking in engines that search for a
Greedy-maximal match directly \citep{cox2007}; the ordering-theoretic
failure identified here is the same failure that shows up operationally
as exponential blowup in those engines. The frontier-based Greedy
parsers of Chapter~\ref{sec:partial-derivative-parser} do not search for
an ordering maximum at all, so this particular failure mode does not
resurface there; Section~\ref{sec:posix-longest-leftmost} returns to
$(a^{*})^{*}$ once POSIX's rule (K1) is available, to show why POSIX
does not have this problem either.

This exposes a limitation of the Greedy ordering: without additional
restrictions or special handling of empty matches, a maximal parse tree
need not exist. The POSIX ordering avoids this problem by treating empty
matches explicitly.

%

\subsection{POSIX Longest-Leftmost Matching}
\label{sec:posix-longest-leftmost}

POSIX uses a different disambiguation policy. Subpatterns are compared
according to the length of the substrings they match, with subpatterns
that occur earlier in the regular expression taking priority. The
standard states that, consistent with the whole match being the
longest of the leftmost matches, each subpattern, from left to right,
shall match the longest possible string \citep{posix_ere}.

Sulzmann and Lu formalize this policy as an ordering on parse trees
\citep{sulzmannlu2014}. Let $\operatorname{len}(w)$ denote the number of
characters in $w$. The ordering is defined by the rules of
Equation~\eqref{eq:posix-order}:

\begin{equation}\label{eq:posix-order}
\begin{array}{ll}
\text{(C1)} &
\dfrac{
  v_1 = v_1'
  \quad
  v_2 \succ_{r_2} v_2'
}{
  (v_1,v_2) \succ_{r_1 r_2} (v_1',v_2')
}
\\[1.2em]

\text{(C2)} &
\dfrac{
  v_1 \succ_{r_1} v_1'
}{
  (v_1,v_2) \succ_{r_1 r_2} (v_1',v_2')
}
\\[1.2em]

\text{(A1)} &
\dfrac{
  \operatorname{len}(\lvert v_2\rvert)
  >
  \operatorname{len}(\lvert v_1\rvert)
}{
  \mathit{Right}\ v_2
  \succ_{r_1+r_2}
  \mathit{Left}\ v_1
}
\\[1.2em]

\text{(A2)} &
\dfrac{
  \operatorname{len}(\lvert v_1\rvert)
  \geq
  \operatorname{len}(\lvert v_2\rvert)
}{
  \mathit{Left}\ v_1
  \succ_{r_1+r_2}
  \mathit{Right}\ v_2
}
\\[1.2em]

\text{(A3)} &
\dfrac{
  v_2 \succ_{r_2} v_2'
}{
  \mathit{Right}\ v_2
  \succ_{r_1+r_2}
  \mathit{Right}\ v_2'
}
\\[1.2em]

\text{(A4)} &
\dfrac{
  v_1 \succ_{r_1} v_1'
}{
  \mathit{Left}\ v_1
  \succ_{r_1+r_2}
  \mathit{Left}\ v_1'
}
\\[1.2em]

\text{(K1)} &
\dfrac{
  \lvert v:vs\rvert = \varepsilon
}{
  [\,]
  \succ_{r^{*}}
  v:vs
}
\\[1.2em]

\text{(K2)} &
\dfrac{
  \lvert v:vs\rvert \neq \varepsilon
}{
  v:vs
  \succ_{r^{*}}
  [\,]
}
\\[1.2em]

\text{(K3)} &
\dfrac{
  v_1 \succ_r v_2
}{
  v_1:vs_1
  \succ_{r^{*}}
  v_2:vs_2
}
\\[1.2em]

\text{(K4)} &
\dfrac{
  v_1 = v_2
  \quad
  vs_1 \succ_{r^{*}} vs_2
}{
  v_1:vs_1
  \succ_{r^{*}}
  v_2:vs_2
}
\end{array}
\end{equation}

The corresponding non-strict relation $\succeq_r$ is defined by

\begin{equation}\label{eq:mvp-33}
v_1 \succeq_r v_2
\quad\Longleftrightarrow\quad
v_1 = v_2
\quad\text{or}\quad
v_1 \succ_r v_2.
\end{equation}

Rules (C1) and (C2) propagate comparisons through concatenation.
Rule (C1) compares the second component when the first components are
equal, while (C2) gives priority to the first component whenever it
differs. Thus, earlier subexpressions have priority over later ones.

Rules (A1)--(A4) define the crucial difference from Greedy matching.
When two alternatives compete, the alternative producing the longer
substring is preferred. If the two substrings have the same length,
(A2) gives preference to the left alternative. Rules (A3) and (A4)
continue the comparison when both candidates have already chosen the
same alternative.

Rules (K1)--(K4) apply the same idea to repetition. In particular,
(K1) prefers the empty list when a non-empty sequence of iterations also
flattens to $\varepsilon$, while (K2) prefers a non-empty match over the
empty list. Rules (K3) and (K4) compare the iterations structurally,
first comparing their heads and then their tails.

%

\paragraph{Applied to $r_1$.}
Return to
$r_1=(a+ab)(b+\varepsilon)$ against $\mathtt{ab}$, with

\begin{equation}\label{eq:mvp-34}
v
=
\bigl(\mathit{Right}(a,b),\mathit{Right}()\bigr)
\end{equation}

and

\begin{equation}\label{eq:mvp-35}
v'
=
\bigl(\mathit{Left}\ a,\mathit{Left}\ b\bigr).
\end{equation}

Both trees are pairs, so the comparison starts at (C2), whose
conclusion $(v_1,v_2) \succ_{r_1 r_2} (v_1',v_2')$ matches with

\begin{equation}\label{eq:mvp-36}
r_1 = a+ab, \quad r_2 = b+\varepsilon, \quad
v_1 = \mathit{Right}(a,b), \quad v_1' = \mathit{Left}\ a .
\end{equation}

(C2) needs only its first components to differ, which they do, so it
reduces the whole comparison to $v_1 \succ_{r_1} v_1'$, that is to
$\mathit{Right}(a,b)$ against $\mathit{Left}\ a$ within $a+ab$. Both
sides are now tagged injections of an alternation, so an $\mathsf{A}$
rule applies, with

\begin{equation}\label{eq:mvp-37}
r_1 = a, \quad r_2 = ab, \quad v_1 = a, \quad v_2 = (a,b) .
\end{equation}

Under that binding (A1)'s premise holds, since the right alternative
consumes strictly more of the input:

\begin{equation}\label{eq:mvp-2}
\begin{aligned}
\operatorname{len}(\lvert v_2\rvert)
  &= \operatorname{len}(\lvert (a,b)\rvert) = 2 \\
\operatorname{len}(\lvert v_1\rvert)
  &= \operatorname{len}(\lvert a\rvert) = 1,
  \qquad\text{so } 2 > 1 .
\end{aligned}
\end{equation}

(A1) is the only rule that can fire here: (A2) would need
$\operatorname{len}(\lvert v_1\rvert) \geq
\operatorname{len}(\lvert v_2\rvert)$, which fails at $1 \geq 2$, and
(A3) and (A4) both require the two sides to carry the \emph{same} tag,
which they do not. So (A1) gives

\begin{equation}\label{eq:mvp-38}
\mathit{Right}(a,b)
\succ_{a+ab}
\mathit{Left}\ a,
\end{equation}

and discharging (C2)'s premise with it gives $v \succeq_{r_1} v'$.
POSIX consequently selects

\begin{equation}\label{eq:mvp-39}
v
=
\bigl(\mathit{Right}(a,b),\mathit{Right}()\bigr),
\end{equation}

which is the opposite of the Greedy result.

%

\paragraph{Applied to $r_2$.}
For
$r_2=(a+b+ab)^{*}$ against $\mathtt{ab}$, consider

\begin{equation}\label{eq:mvp-40}
v_2 =
[\mathit{Right}\ \mathit{Right}(a,b)]
\end{equation}

and

\begin{equation}\label{eq:mvp-41}
v_2' =
[\mathit{Left}\ a,\mathit{Right}\ \mathit{Left}\ b].
\end{equation}

Both trees are non-empty iteration lists, so (K3) applies. Matching its
conclusion, which is $v_1{:}vs_1 \succ_{r^{*}} v_2{:}vs_2$, gives

\begin{equation}\label{eq:mvp-42}
r = a+(b+ab), \quad
v_1 = \mathit{Right}\ \mathit{Right}(a,b), \quad
v_2 = \mathit{Left}\ a,
\end{equation}

and $vs_1 = [\,]$, $vs_2 = [\mathit{Right}\ \mathit{Left}\ b]$. (K3)
compares only the first iterations and ignores the tails entirely, which
is what reduces the question to $v_1 \succ_r v_2$ inside the star body.
Within that body both sides are tagged injections, binding

\begin{equation}\label{eq:mvp-43}
r_1 = a, \quad r_2 = b+ab, \quad
v_1 = a, \quad v_2 = \mathit{Right}(a,b),
\end{equation}

and (A1) fires again because the right alternative consumes more:

\begin{equation}\label{eq:mvp-44}
\operatorname{len}(\lvert\mathit{Right}(a,b)\rvert)
=
2
>
1
=
\operatorname{len}(\lvert a\rvert).
\end{equation}

Thus

\begin{equation}\label{eq:mvp-45}
\mathit{Right}\ \mathit{Right}(a,b)
\succ_{a+(b+ab)}
\mathit{Left}\ a,
\end{equation}

and discharging (K3)'s premise with that lifts the comparison to the
star itself:

\begin{equation}\label{eq:mvp-46}
v_2
\succeq_{r_2}
v_2'.
\end{equation}

POSIX therefore selects the single two-character iteration, whereas
Greedy selects the two one-character iterations.

%

\paragraph{Applied to $r_3$.}
For $r_3=(a+a)^{*}$ against $\mathtt{aa}$, the star body binds
$r_1 = a$ and $r_2 = a$, so both alternatives are the same literal and
every candidate at every position matches exactly one character.
(A1)'s premise $\operatorname{len}(\lvert v_2\rvert) >
\operatorname{len}(\lvert v_1\rvert)$ therefore reads $1 > 1$ and never
holds, so every comparison falls through to (A2), whose premise is the
non-strict one:

\begin{equation}\label{eq:mvp-47}
\operatorname{len}(\lvert a\rvert)
\geq
\operatorname{len}(\lvert a\rvert)
\quad\Longrightarrow\quad
\mathit{Left}\ a \succ_{a+a} \mathit{Right}\ a.
\end{equation}

(K3) and (K4) lift this position by position, giving the same result as
Greedy:

\begin{equation}\label{eq:mvp-48}
[\mathit{Left}\ a,\mathit{Left}\ a].
\end{equation}

Unlike $r_1$ and $r_2$, Greedy and POSIX agree here. This is not a
coincidence: (A2) only ever overrules (A1) when both alternatives tie in
length, and a tie is exactly the condition under which Greedy's
unconditional left-preference and POSIX's tie-broken left-preference
compute the same thing, the same connection already observed for
\texttt{"if"} in Section~\ref{sec:matching-problem}.

%

\paragraph{Applied to $O$.}
For $O=((A+B)+C)+D$ against $w=\mathtt{239}$, only $A$ and $C$ match
at all, so the comparison is between a tree choosing $A$ and a tree
choosing $C$. Both trees agree on the outer alternation: neither
chooses $D$, so the comparison reduces, by (A4), to a comparison
within $(A+B)+C$ of a tree choosing $A$ (further inside $A+B$) against
one choosing $C$ directly. Since $A$ and $C$ both consume the full
input,

\begin{equation}\label{eq:mvp-49}
\operatorname{len}(\lvert v_A\rvert)
=
3
=
\operatorname{len}(\lvert v_C\rvert),
\end{equation}

so (A1)'s strict inequality fails and (A2) resolves the tie in favor
of the left-hand candidate, the one going through $A+B$ and ultimately
choosing $A$. POSIX therefore also selects $A$. As with $r_3$, this is
not a coincidence: $A$ and $C$ tie in length, so (A2)'s tie-break
coincides with Greedy's unconditional left-preference, even though $A$
and $C$, unlike $r_3$'s two copies of \texttt{a}, are not otherwise
interchangeable.

%

\paragraph{Applied to the infinite-chain examples.}
Rule (K1) resolves Section~\ref{sec:infinite-empty-chain}'s pathology
directly: for $\varepsilon^{*}$ against $\varepsilon$, since

\begin{equation}\label{eq:mvp-50}
\lvert () : [\,]\rvert
=
\lvert ()\rvert\,\lvert [\,]\rvert
=
\varepsilon,
\end{equation}

(K1) gives

\begin{equation}\label{eq:mvp-51}
[\,]
\succ_{\varepsilon^{*}}
[()].
\end{equation}

The same rule applies to every longer list of empty iterations. Thus
$[\,]$ is preferred to every parse tree obtained by repeatedly matching
$\varepsilon$.

For $(\varepsilon+a)^{*}$ against $\mathtt{a}$, the relevant comparison
is different. Consider

\begin{equation}\label{eq:mvp-52}
v_0=[\mathit{Right}\ a]
\end{equation}

and

\begin{equation}\label{eq:mvp-53}
v_1=[\mathit{Left}\ (),\mathit{Right}\ a].
\end{equation}

Rule (K3) compares their first iterations, reducing the comparison to

\begin{equation}\label{eq:mvp-54}
\mathit{Right}\ a
\quad\text{versus}\quad
\mathit{Left}\ ()
\end{equation}

within $\varepsilon+a$. Rule (A1) applies because

\begin{equation}\label{eq:mvp-55}
\operatorname{len}(\lvert a\rvert)
=
1
>
0
=
\operatorname{len}(\lvert ()\rvert).
\end{equation}

Therefore

\begin{equation}\label{eq:mvp-56}
\mathit{Right}\ a
\succ_{\varepsilon+a}
\mathit{Left}\ (),
\end{equation}

and hence

\begin{equation}\label{eq:mvp-57}
v_0
\succ_{(\varepsilon+a)^{*}}
v_1.
\end{equation}

The same comparison applies to every parse tree obtained by inserting
additional empty iterations before the iteration matching
$\mathtt{a}$. Thus the POSIX ordering selects $v_0$, while the Greedy
ordering admits the infinite ascending chain.

Finally, for $(a^{*})^{*}$ against $\mathtt{aaa}$, compare

\begin{equation}\label{eq:mvp-58}
[[a,a,a]]
\qquad\text{against}\qquad
[[a,a,a],[\,]].
\end{equation}

Since $\lvert [\,]:[\,]\rvert = \varepsilon$, (K1) applies to the tail
once the first (identical) iteration is matched by (K4), giving

\begin{equation}\label{eq:mvp-59}
[[a,a,a]]
\succ_{(a^{*})^{*}}
[[a,a,a],[\,]],
\end{equation}

and the same argument disposes of every list padded with further empty
iterations after the first. POSIX therefore selects the single outer
iteration whose inner star consumes all three characters directly,

\begin{equation}\label{eq:mvp-60}
[[a,a,a]],
\end{equation}

exactly the same tree the running implementation returns for
\texttt{(a*)*} against \texttt{aaa} under \texttt{deriv\_std\_rec} and
\texttt{deriv\_bc}. Unlike the abstract Greedy ordering, which has no
maximal element for this expression, the frontier-based Greedy
implementations of Chapter~\ref{sec:partial-derivative-parser} do not
compute an ordering maximum at all, and return this same tree without
difficulty; the ordering-theoretic failure identified in
Section~\ref{sec:infinite-empty-chain} is a property of the abstract
Greedy comparison, not of every construction that happens to compute a
Greedy answer.

The contrast between the two policies is now explicit. Greedy resolves
ambiguity according to the structural priority encoded by the regular
expression, whereas POSIX compares the substrings assigned to
subexpressions and gives priority to longer leftmost matches.
Sulzmann and Lu prove that, unlike the Greedy ordering, the POSIX
ordering is total and has a maximal element for every expression and
input \citep{sulzmannlu2014}.
The derivative parser developed in Chapter~\ref{sec:derivative-parser} computes that
candidate without enumerating all possible parse trees.